\section{Bruun's Real Factor Tree}
\label{sec:bruun}

Bruun's construction begins with the observation that a real DFT need not split
$z^N-1$ into complex linear factors.  The two real roots are represented by
$z-1$ and $z+1$; every other conjugate pair is represented by a real quadratic.
The sparse route to those factors is a repeated half-angle identity.

\subsection{Binomial spine and trinomial descendants}

For every positive integer $m$ define
\begin{equation}
 Q_{m,\theta}(z)=z^{2m}-2\cos\theta\,z^m+1.
 \label{eq:Q-definition}
\end{equation}
The coefficient field throughout the exact development is
\begin{equation}
\K_N=\mathbb{Q}\bigl(\sqrt2,
\{\cos(2\pi k/N),\sin(2\pi k/N):k\in\Z\}\bigr)
 \subset\R.
 \label{eq:coefficient-field}
\end{equation}
All polynomial quotients in Sections \ref{sec:bruun}--\ref{sec:fourier-identity}
may be taken over $\K_N$.

\begin{lemma}[Half-angle factorization]
For even $m$,
\begin{equation}
 Q_{m,\theta}
 =Q_{m/2,\theta/2}\,Q_{m/2,\pi-\theta/2}.
 \label{eq:half-angle-factorization}
\end{equation}
\end{lemma}
\begin{proof}
Put $v=z^{m/2}$ and $a=2\cos(\theta/2)$.  The right side is
$(v^2-av+1)(v^2+av+1)=(v^2+1)^2-a^2v^2$.
Since $a^2=2+2\cos\theta$, this equals
$v^4-2\cos\theta\,v^2+1$.
\end{proof}

The root polynomial first splits along the binomial spine,
\begin{equation}
 z^{2m}-1=(z^m-1)(z^m+1),
 \label{eq:binomial-spine}
\end{equation}
and $z^m+1=Q_{m/2,\pi/2}$.  Repeated application of
\eqref{eq:half-angle-factorization} terminates at
\begin{equation}
 q_k(z)=z^2-2\cos(2\pi k/N)z+1,
 \quad 1\leq k<N/2,
 \label{eq:terminal-quadratic}
\end{equation}
together with $z-1$ and $z+1$.  Each $q_k$ represents the conjugate frequencies
$k$ and $N-k$.

\subsection{Forward reduction}

If $p=p_-p_+$ with coprime children, the local forward map is simply
\begin{equation}
 [f]_p\longmapsto([f]_{p_-},[f]_{p_+}).
 \label{eq:local-reduction}
\end{equation}
Composing \eqref{eq:local-reduction} down the complete factor tree gives
\begin{equation}
 \red_N:R_N\longrightarrow
 \R\times\R\times
 \prod_{k=1}^{N/2-1}\R[z]/(q_k).
 \label{eq:global-reduction}
\end{equation}
The two scalar components correspond to $z=1$ and $z=-1$.

\subsection{The resumed questions}

The line of inquiry recorded for the present construction began in November
2024 with four concrete problems: reproduce Bruun's reduction; identify the
leaf shuffle; construct an inverse; and determine whether a normalized form
could serve as a complete transform.  On June 10, 2026 the thread was resumed
from the inverse problem.  Claude Fable contributed the decisive local
Chinese-remainder reconstruction described in Section \ref{sec:crt}.  That
formula made the factor tree bidirectional and exposed the next obstruction:
exact inversion alone does not choose a well-conditioned coordinate system.

The subsequent questions are stated here because they determine the logical
shape of the paper.  First, can the terminal quotient be given a unit complex
coordinate without introducing complex polynomial arithmetic?  Second, do
alternative balanced bases improve the local factor geometry?  Third, which
orders preserve the tree's sparse intervals?  Fourth, do downward, upward, and
diagonal walks compute the same boundary map?  The answers will be a normalized
quadratic chart, a covering-map locality principle, an address algebra, and
three exact schedules of one Fourier factorization.

\subsection{Dated reconstruction of the inquiry}

Table~\ref{tab:development-line} records the order in which the questions were
resolved.  The logical order of the proofs is different: normalization must be
established before any schedule can inherit its metric.  The June 14--16
Chebyshev work was a return to alternative bases after the first DIT
construction, not the origin of the normalized chart present on June 10.

\begin{table}[t]
\caption{Documented mathematical development leading to the present theory.}
\label{tab:development-line}
\centering
\footnotesize
\begin{tabular}{@{}p{0.20\columnwidth}p{0.73\columnwidth}@{}}
\toprule
Date & Mathematical event\\
\midrule
Nov. 2024 & Bruun reduction, order, inverse, and viability posed as questions.\\
June 10, 2026 & Fable CRT inverse; first staged normalized forward and inverse maps.\\
June 10--11 & Equal-angle block and real rotation inquiry.\\
June 12 & Subtree locality, output order, and first DIT construction.\\
June 13 & Faithful monomial-residue reference reconstructed.\\
June 14--16 & Radix-four and Chebyshev alternative-basis study.\\
June 23 & Exact odd-frequency transform defined independently.\\
July 4--5 & Phase-packet shear, DIP, and twisted span walk.\\
July 5 & Positive-cone and field-delay realizations.\\
Aug. 13 & Corrected cone composition and mechanical gauges.\\
\bottomrule
\end{tabular}
\end{table}

\begin{figure}[t]
\centering
\begin{tikzpicture}[level distance=8mm,
 level 1/.style={sibling distance=38mm},
 level 2/.style={sibling distance=21mm},
 level 3/.style={sibling distance=12mm},
 every node/.style={font=\tiny,inner sep=1pt},
 edge from parent/.style={draw,-{Latex[length=1.2mm]}}]
\node {$z^8-1$}
 child {node {$z^4-1$}
   child {node {$z^2-1$}
     child {node {$z-1$}}
     child {node {$z+1$}}}
   child {node {$z^2+1$}}}
 child {node {$z^4+1$}
   child {node {$z^2-\sqrt2z+1$}}
   child {node {$z^2+\sqrt2z+1$}}};
\end{tikzpicture}
\caption{The $N=8$ real factor tree.  The leaves represent DC, Nyquist, and
three conjugate frequency pairs.  Every edge is a remainder map; Section
\ref{sec:crt} reverses every edge explicitly.}
\label{fig:n8-tree}
\end{figure}
